% ==================== CHAPTER 2: LITERATURE REVIEW ====================
\section{Literature Review}

\subsection{EGFR as a Therapeutic Target in NSCLC}

The Epidermal Growth Factor Receptor (EGFR), also known as ErbB1 or HER1, is a transmembrane glycoprotein and a member of the ErbB family of receptor tyrosine kinases (RTKs) \parencite{Siegeletal2023}. Structurally, EGFR comprises an extracellular ligand-binding domain, a single transmembrane helix, an intracellular juxtamembrane region, a tyrosine kinase domain, and a C-terminal regulatory tail. Under physiological conditions, ligand binding (e.g., epidermal growth factor [EGF] or transforming growth factor-alpha [TGF-$\alpha$]) induces receptor dimerization, autophosphorylation of specific tyrosine residues, and subsequent activation of multiple downstream signaling pathways. The primary pathways include the RAS-RAF-MEK-ERK (MAPK) pathway, which promotes cell proliferation, and the PI3K-AKT-mTOR pathway, which enhances cell survival and inhibits apoptosis \parencite{Siegeletal2023, RotowJanne2023}.

In non-small cell lung cancer (NSCLC), activating somatic mutations in the EGFR kinase domain lead to ligand-independent constitutive signaling, driving oncogenesis. The two most prevalent sensitizing mutations are in-frame deletions in exon 19 (Ex19del, approximately 45\% of cases) and the L858R point mutation in exon 21 (approximately 40\% of cases). Together, these mutations account for 85--90\% of EGFR-mutant NSCLC \parencite{RotowJanne2023}. The overall prevalence of EGFR mutations in NSCLC varies by ethnicity: 10--15\% in Caucasian populations and 40--50\% in East Asian populations \parencite{Siegeletal2023}. This ethnic disparity is particularly relevant in regions such as Nepal and South Asia, where NSCLC is often diagnosed at advanced stages and access to molecular testing and targeted therapies remains limited. Consequently, EGFR has become one of the most important therapeutic targets in precision oncology for NSCLC, with EGFR tyrosine kinase inhibitors (TKIs) transforming patient outcomes in mutation-positive cases \parencite{Passaroetal2021}.

\subsection{Generations of EGFR TKIs and Resistance Mechanisms}

The evolution of EGFR TKIs reflects an iterative response to emerging resistance mechanisms, progressing through three distinct generations \parencite{RotowJanne2023}.

First-generation TKIs, such as gefitinib and erlotinib, are reversible, ATP-competitive inhibitors that bind preferentially to mutant EGFR. These agents produced unprecedented response rates and progression-free survival benefits in patients with sensitizing mutations. However, acquired resistance developed in nearly all patients within 9--13 months, predominantly through the secondary T790M ``gatekeeper'' mutation in exon 20. This substitution increases ATP affinity while sterically hindering drug binding \parencite{Passaroetal2021}.

Second-generation TKIs, including afatinib and dacomitinib, are irreversible covalent inhibitors targeting Cys797 in the ATP-binding pocket. They exhibit broader activity against T790M in preclinical models and inhibit other ErbB family members. Despite improved potency, their clinical use has been restricted by on-target toxicity to wild-type EGFR, leading to dose-limiting adverse effects such as skin rash and diarrhea \parencite{Passaroetal2021, RotowJanne2023}.

Third-generation TKIs, most notably osimertinib, were rationally designed to covalently bind T790M-mutant EGFR while sparing wild-type EGFR. Osimertinib is currently the preferred first-line therapy for EGFR-mutant NSCLC, demonstrating superior efficacy and central nervous system penetration compared to earlier generations. Nevertheless, resistance inevitably emerges, most commonly via the C797S mutation at the covalent binding site. This leads to triple-mutant EGFR (e.g., Ex19del/T790M/C797S or L858R/T790M/C797S), which is resistant to all approved TKIs and represents a major unmet clinical need in advanced NSCLC treatment \parencite{RotowJanne2023, Maliketal2025}. 

To address this critical gap, fourth-generation EGFR TKIs, such as BLU-945, are currently undergoing clinical evaluation. These novel agents are specifically designed as reversible, wild-type-sparing inhibitors that maintain robust \textit{in vivo} potency against the triple-mutant configurations that mediate third-generation resistance \parencite{Limetal2024}.

\subsection{Bioactivity Data and the ChEMBL Database}

The ChEMBL database, maintained by the European Molecular Biology Laboratory-European Bioinformatics Institute (EMBL-EBI), is one of the largest and most widely utilised public repositories of bioactive molecules. Following recent comprehensive updates, ChEMBL now contains over 2.3 million distinct compounds and more than 20 million bioactivity measurements across thousands of biological targets \parencite{Zdraziletal2024}. For EGFR research, the relevant entry is target ID CHEMBL203, which aggregates bioactivity data from diverse published and deposited sources \parencite{Zdraziletal2024}.

A major challenge when utilising ChEMBL data for machine learning is the inherent heterogeneity of experimental conditions. Bioactivity values (primarily IC$_{50}$) originate from hundreds of different laboratories employing varied assay protocols, cell lines, and conditions. To mitigate this, rigorous, modern curation practices mandate strict filtering for exact IC$_{50}$ values, conversion to the logarithmic pIC$_{50}$ scale, removal of ambiguous data points, and aggregation of multiple measurements per compound \parencite{Huangetal2021}. In the present study, additional standardisation using RDKit (salt removal, canonicalisation, and molecular weight filtering between 100--1000 Da) yielded a high-quality dataset of 10,523 unique compounds, with approximately 74\% classified as active at the commonly used threshold of pIC$_{50} \geq 6.0$ (equivalent to IC$_{50} \leq 1~\mu$M) \parencite{Huangetal2021, Maliketal2025}.

\subsection{Molecular Representation and Traditional Machine Learning Approaches}

Prior to deep learning approaches, molecular property prediction relied heavily on fixed-length molecular fingerprints paired with classical machine learning algorithms. The Extended Connectivity Fingerprint (ECFP), which encodes local atomic environments by iteratively capturing circular neighbourhoods, has long served as a de facto baseline standard in computational cheminformatics \parencite{Dengetal2022}. ECFP4 (typically folded to 2048 bits) provides a balance of computational efficiency and chemical informativeness.

These fingerprints are commonly combined with ensemble methods such as Random Forest (RF) or simpler linear models like Logistic Regression. Random Forest constructs multiple decision trees on bootstrapped data subsets and aggregates predictions, providing robustness to noise and intrinsic feature importance rankings \parencite{Jiangetal2021}. Logistic Regression, while less complex, offers high interpretability through coefficient analysis. Despite their widespread adoption and strong baseline performance, these approaches suffer from a fundamental limitation: the molecular representation is hand-crafted and static, potentially missing higher-order topological patterns and long-range atomic interactions critical to biological activity \parencite{Jiangetal2021, Dengetal2022}.

\subsection{Graph Neural Networks for Molecular Property Prediction}

Graph Neural Networks (GNNs) have emerged as a transformative approach in molecular machine learning by learning representations directly from the molecular graph structure rather than relying on pre-defined fingerprints \parencite{Gilmeretal2017}. In the message-passing framework, each atom (node) iteratively aggregates features from its neighbouring atoms (edges) across multiple layers. This process enables the model to encode increasingly complex chemical contexts. After the message-passing phase, a readout function produces a fixed-size graph-level embedding suitable for classification tasks \parencite{Gilmeretal2017}.

Prominent GNN architectures evaluated in recent drug discovery frameworks include the Graph Convolutional Network (GCN), which performs degree-normalised aggregation, and the Graph Attention Network (GAT), which introduces learnable attention coefficients \parencite{Zhangetal2021}. More advanced architectures, such as AttentiveFP, employ multi-step attentive pooling and natively incorporate edge features, demonstrating particular strength in capturing local and global pharmacophore information through hierarchical attention mechanisms \parencite{Xiongetal2020, Jiangetal2021}.

Large-scale studies have confirmed that GNNs frequently outperform traditional fingerprint-based methods on structurally diverse datasets, and benefit significantly from self-supervised pre-training on millions of unlabelled molecules \parencite{Huetal2020}. Recently, the field has also seen a rapid expansion into 3D geometric deep learning. Geometric GNNs explicitly incorporate three-dimensional atomic coordinates and spatial symmetries (e.g., rotational equivariance), offering more expressive representations of protein-ligand binding conformations than standard 2D topological graphs \parencite{Atzetal2021}.

\subsection{GNN Models for EGFR Inhibitor Prediction}

One of the most relevant prior works is DeepEGFR, a GAT-based model trained on ChEMBL-derived EGFR bioactivity data \parencite{Maliketal2025}. DeepEGFR achieved strong classification performance in distinguishing active from inactive EGFR inhibitors. However, several limitations remain: reliance on random data splitting (which often overestimates generalisation to novel scaffolds), absence of instance-level structural explanations, and lack of prospective application to virtual screening for new lead compounds \parencite{Maliketal2025, Jiangetal2021}.

The present thesis extends this line of research by systematically evaluating four GNN architectures (GCN, GAT, GIN, and AttentiveFP), selecting AttentiveFP as the primary model for its advanced attention mechanisms and edge feature integration. It employs a rigorous Bemis-Murcko scaffold split (80/10/10), utilises comprehensive 29-dimensional node and 12-dimensional edge features, and conducts Bayesian hyperparameter optimisation using Optuna, thereby addressing key methodological gaps in earlier EGFR GNN studies \parencite{Jiangetal2021}.

\subsection{Explainability in GNN-Based Drug Discovery}

Despite their predictive power, deep learning models in drug discovery have historically been criticised as ``black boxes,'' limiting their acceptance by medicinal chemists in practical settings \parencite{Luoetal2022}. Post-hoc explainability techniques have been developed to bridge this gap. GNNExplainer \parencite{Yingetal2019} is a model-agnostic method that identifies the minimal subgraph (through node and edge masks) most responsible for a given prediction by maximising mutual information between the masked input and the model output. 

Applied to EGFR inhibitors such as gefitinib, erlotinib, osimertinib, and lapatinib, GNNExplainer can highlight pharmacophoric features (e.g., hinge-binding nitrogens or hydrophobic halogens) in a chemically interpretable manner \parencite{Yingetal2019, Luoetal2022}. While tools such as SHAP provide global feature importance, the current work focuses on GNNExplainer for instance-level insights that can be directly compared against known EGFR kinase pharmacophore models.

\subsection{Virtual Screening with Machine Learning}

Virtual screening (VS) is a cornerstone of computational drug discovery, enabling the prioritisation of large chemical libraries for experimental validation \parencite{Wangetal2023, Sunetal2021}. In ligand-based machine learning VS, a model trained on known active and inactive compounds is deployed to score external libraries, ranking molecules by predicted probability of activity.

This thesis implements a focused VS pipeline by curating compounds active against related kinases (VEGFR2, BRAF, CDK2, ABL1, MET) from ChEMBL while strictly excluding the EGFR training set. After drug-likeness filtering (MW 250--600 Da, LogP -1 to 7), the optimised AttentiveFP model scores 1,872 unseen compounds. Top candidates are further prioritised using RDKit-computed ADMET properties, including Lipinski compliance and the Quantitative Estimate of Drug-likeness (QED) score \parencite{Wangetal2023, Tropshaetal2020}.

Notable successes, such as the discovery of novel antibiotics through GNN-based screening of over 100 million compounds \parencite{Stokesetal2020}, highlight the powerful potential of this approach. However, most published EGFR GNN models have stopped at classification without extending to prospective lead identification \parencite{Stokesetal2020, Maliketal2025}.

\subsection{Summary and Research Gap}

The reviewed literature demonstrates substantial progress in EGFR-targeted therapy, bioactivity data curation, traditional cheminformatics, and GNN-based modelling. Nevertheless, critical gaps persist. Existing GNN studies for EGFR inhibitors typically lack one or more of the following: rigorous scaffold-based evaluation, advanced architectures with rich node/edge features, systematic interpretability using GNNExplainer on clinical compounds, and integration with targeted virtual screening pipelines for novel lead discovery \parencite{Maliketal2025, Jiangetal2021, Luoetal2022}.

This thesis addresses these gaps through a comprehensive, end-to-end computational pipeline that combines high-quality ChEMBL curation, scaffold-split AttentiveFP modelling with Bayesian optimisation, GNNExplainer-based interpretability aligned with EGFR pharmacophores, and virtual screening of related-kinase inhibitors. By doing so, it contributes both methodological advancements in molecular machine learning and practical candidates for addressing drug-resistant NSCLC.

\newpage